% Drawing 74/04-7 -- folder 74, batch 04, archive page 73 (PDF page 74).
% Transcription: un hexagone aux sommets successifs $a$, $b'$, $c$, $a'$, $b$, $c'$.
% Kind: hexagon with labelled vertices. Method 3 (encoded, rendered from the data). Status: ok.
% Notes: « Soit h = » a hexagon with successive vertices a, b', c, a', b, c' (the labels a and c are faint). Encoded as polyhedron/1 (one 2-cell, the hexagon as a vertex cycle).
% Coordinates: 1 unit = 1 mm = 1/100 of the crop width (bbox in page fractions: [0.18, 0.09, 0.45, 0.2]).
% Redrawn by Claude (claude-opus-5-5), issue #73, 2026-11-07; not checked by a human; not compiled (no TeX in the session).
% Rights: reproduces part of the work -- behind RENDER_TRANSCRIPTS, never published by itself.
\documentclass[tikz,border=3pt]{standalone}
\usepackage{amsmath,amssymb}
\usetikzlibrary{patterns,arrows.meta}
\tikzset{pencil/.style={black!45}, struck/.style={black!35}, mark/.style={circle,draw,inner sep=1.3pt}, pt/.style={circle,fill,inner sep=0.9pt}, opt/.style={circle,draw,fill=white,inner sep=1.1pt}}
\begin{document}
\begin{tikzpicture}[x=1mm,y=1mm,line cap=round,line join=round,font=\scriptsize]
\draw (60,-22) -- (79,-21) -- (88,-34) -- (80,-46) -- (66,-46) -- (57,-35) -- cycle;
\node[inner sep=1pt] at (59,-18) {\scriptsize $a$};
\node[inner sep=1pt] at (83,-18) {\scriptsize $b'$};
\node[inner sep=1pt] at (92,-34) {\scriptsize $c$};
\node[inner sep=1pt] at (85,-49) {\scriptsize $a'$};
\node[inner sep=1pt] at (66,-51) {\scriptsize $b$};
\node[inner sep=1pt] at (52,-35) {\scriptsize $c'$};
\end{tikzpicture}
\end{document}
